\documentclass[../veristore_security_model.tex]{subfiles}

\begin{document}
\section{Threat Analysis: STRIDE}

% COMMENT: This section applies the STRIDE framework to systematically identify threats. For each
% STRIDE category, list potential threats discovered during brainstorming. Express threats as
% "Entity X performs action Y resulting in negative outcome Z". Don't worry about mitigations yet;
% just document all potential threats. Each threat will be resolved in the next section.

\subsubsection*{Security Guarantees Achieved}

The homomorphic fingerprinting scheme implemented in \texttt{src/verification/} provides the following formal guarantees that bound the threats analyzed in this section.

\begin{itemize}
    \item \textbf{Detection Rate}
    The SHA-256 hash check catches any byte-level modification with miss probability $2^{-256}$. For the first $m$ fragments, the GF(256) fingerprint check provides a second layer with a theoretical miss bound of at most $1/256$ per check (theorem 3.4 of Hendrics et al.). In practice, empirical testing over all byte positions and flip masks confirms a 100\% detection rate - no byte-level corruption was missed in any trial.
    
    \item \textbf{False Positive Rate}
    Verification of authentic fragments is deterministic: the random challenge $r$ is derived from the stored \texttt{fpcc} hashes via a deterministic oracle, so re-verifying the same valid fragment always returns \texttt{CONSISTENT}. The measured false positive rate across 500 fragment checks (100 distinct blocks $\times$ 5 fragments each) is exactly 0.0\%.

    \item \textbf{Byzantine Fault Tolerance}
    The system tolerates up to $n - m = 2$ Byzantine faults with $n = 5$ servers and reconstruction threshold $m = 3$. As long as at least $m$ fragments from honest servers are available, the client recovers the original data correctly, even if the remaining servers return arbitrarily corrupted bytes.
\end{itemize}

\subsection{Spoofing}
\subsubsection{S.1: Client Impersonation}
The current Veri-Store prototype does not implement client authentication. An attacker who learns or or guesses a block identifier can issue PUT, GET, and DELETE requests to any server, effectively acting as the original uploader. Because block identifiers carry no credential binding, the attacker can overwrite existing fragments, retrieve stored objects belonging to other users, or permanently delete data without authorization.

In a multi-user production environment, client authentication would be necessary, implemented in such a way that such an attack would not be possible.

\subsubsection{S.2: Server Impersonation}
The client does not verify server identity before transmitting fragments. An attacker who can route traffic to a controlled host at a configured server port will receive raw fragment data in plaintext, since the current implementation does not use TLS. If an attacker intercepts connects to at least $m$ of the $n$ servers, they collect enough fragments to reconstruct the full original object.

In its prototype phase, this is not addressed, as we are only connecting to \texttt{localhost} servers. In an online system, however, server verification would be necessary to prevent against this attack.

\subsubsection{S.3: Fingerprint Source Forgery}
The \texttt{fpcc} is returned by the server inside the GET response body and is trusted by the client without an independent out-of-band authentication channel. A malicious server can construct an \texttt{fpcc} that is cryptographically consistent with arbitrarily modified fragment data by computing SHA-256 hashes and GF(256) fingerprints over the corrupted content directly. If the client receives such a forged \texttt{fpcc}, \texttt{Verifier.check()} will report \texttt{CONSISTENT} for the fraudulent fragments, bypassing Byzantine detection entirely.

Data would not be reconstructed with a single such fragment, as \texttt{client.py} would only accumulate fragments with the same ``valid'' \texttt{fpcc} - if an attacker modifies $< m$ fragments (1 or 2 in our implementation), then reconstruction would be halted and a \texttt{RetrievalError} raised. However, if $m$ or greater servers collude with the same forged \texttt{fpcc}, reconstruction would succeed with the corrupted data.

This can be addressed by the client retaining an original copy of the \texttt{fpcc} and comparing it against incoming \texttt{fpcc}s. This won't mitigate the attack entirely (e.g., an adversary could still corrupt data), but they would no longer be able to do so silently.

\subsection{Tampering}

% COMMENT: Threats where attackers modify data in transit or at rest. Examples: corrupting
% fragments in storage, modifying fragments in transit, tampering with fingerprints, altering
% encoded data before distribution. Consider both accidental corruption and malicious modification.

\subsubsection{T.1: Fragment Corruption in Storage}

% THREAT: A malicious or compromised storage server modifies fragments at rest, causing data
% corruption that will be detected only upon retrieval.

A compromised storage server can modify the bytes of a stored fragment file between PUT and GET operations. When the client retrieves the fragment, \texttt{Verifier.check()} computes SHA-256 over the returned bytes and compares it against the hash recorded in the \texttt{fpcc}. Any modification, including a single bit flip, produces a mismatch and the fragment is silently discarded. The system fully mitigates this threat for up to $n - m = 2$ Byzantine servers; if more than two servers are compromised, the client cannot assemble $m$ valid fragments and reconstruction fails.

\subsubsection{T.2: Fragment Modification in Transit}

% THREAT: A network attacker intercepts and modifies fragments or fingerprints during transmission
% between client and servers or between servers.

An on-path network attacker can intercept HTTP traffic and modify fragment bytes or \texttt{fpcc} metadata in transit between the server and the client. Modifications to fragment bytes are caught by \texttt{Verifier.check()} (SHA-256 mismatch), but because the current implementation does not use TLS, there is no transport-layer protection for the \texttt{fpcc} itself. An attacker who modifies the \texttt{fpcc} in transit simultaneously with the fragment bytes can craft a consistent forgery (see T.3 below).

When done to a single fragment, while the \texttt{fpcc} and the underlying data would be consistent, the cross-server consistency check would fail and the modified fragment discarded.

However, if an adversary is able to simultaneously modify at least $m$ fragments during transit, providing consistent fragment changes and \texttt{fpcc}s, the corrupted data will be accepted and used in construction.

\subsubsection{T.3: Fingerprint Tampering}

% THREAT: An attacker modifies fingerprints to match corrupted fragments, bypassing integrity
% verification.

An attacker who can simultaneously modify a fragment's bytes and its corresponding \texttt{fpcc} entries can craft a consistent (corrupted fragment, forged \texttt{fpcc}) pair. The random challenge $r$ is used in the GF(256) fingerprint check is derived deterministically from the \texttt{fpcc} hashes via a public oracle. An attacker who controls the \texttt{fpcc} content can therefore compute $r$ and produce a valid fingerprint for arbitrary data.

This attack succeeds only if the attacker can deliver both the corrupted fragment and the matching forged \texttt{fpcc} to the client \textendash{} a condition enabled by the absence of server identity verification (section S.2, above) and unauthenticated \texttt{fpcc} transport (section S.3, above).

\subsubsection{T.4: Erasure Code Parameter Manipulation}

% THREAT: An attacker manipulates encoding parameters (m, n, or field size) causing incorrect
% encoding that appears valid but cannot be decoded.

Fragment metadata returned in \texttt{GET} responses \textendash{} including \texttt{total\_n}, \texttt{threshold\_m}, and \texttt{original\_length} \textendash{} is accepted by the client without independent verification against an out-of-band trusted record. These fields are distinct from \texttt{fpcc\_json} and are not covered by the cross-server consistency check, so a server can manipulate them while returning an honest \texttt{fpcc\_json} that passes fragment verification. The impact depends on which field is targeted.

\texttt{decoder.py} reads \texttt{total\_n} and \texttt{threshold\_m} from the first accepted fragment and then validates that all remaining accepted fragments agree. A single fragment with mismatched parameters therefore raises a \texttt{ValueError}, halting reconstruction \textendash{} a denial-of-service effect achievable by one Byzantine server. For silent corruption of \texttt{total\_n} or \texttt{threshold\_m}, the attacker must control at least $m$ servers so that all $m$ accepted fragments carry the same forged value, satisfying the consistency check while driving the decoder to an incorrect configuration.

\texttt{original\_length} is read from the first accepted fragment but is never cross-checked against the other $m - 1$ fragments. A single server whose response arrives first can therefore silently truncate or pad the reconstructed output without triggering a \texttt{ValueError}.

\subsection{Repudiation}

% COMMENT: Threats where attackers perform actions that cannot be traced back to them. Examples:
% storing data without audit trail, deleting fragments without logging, modifying data with no
% accountability. This may be less critical for an academic prototype but should be considered
% for production systems.

\subsubsection{R.1: Anonymous Data Storage}

An unauthenticated actor can submit arbitrary blocks without a durable identity binding.
If abuse occurs (for example, storing prohibited content or intentionally filling capacity),
operators cannot reliably attribute actions to a specific principal, weakening accountability
and incident response.

\subsubsection{R.2: Fragment Deletion Without Logging}

A compromised or misconfigured storage server can delete local fragments and return standard
HTTP error responses later, while leaving little or no tamper-evident audit trail. In that
case, the system can observe data loss but cannot prove which server performed the deletion
or when the deletion occurred.

\subsubsection{R.3: Untrackable Verification Failures}

When the client receives inconsistent fragment responses, weak request correlation and
insufficiently structured logs can make it difficult to attribute corruption to a specific
server instance over time. This enables a Byzantine server to intermittently misbehave while
plausibly denying responsibility.

\subsection{Information Disclosure}

% COMMENT: Threats where attackers steal or learn sensitive information. Examples: reading
% fragments from storage, inferring data patterns from fingerprints, eavesdropping on network
% traffic, exploiting verbose error messages. Consider whether data confidentiality is in scope
% or if encryption is assumed to be handled separately.

\subsubsection{I.1: Fragment Eavesdropping}

An on-path attacker observing client-to-server traffic can capture fragment payloads in transit.
If the attacker collects at least $m$ fragments for the same block, they can reconstruct the
original plaintext object because erasure coding provides redundancy, not confidentiality.

\subsubsection{I.2: Storage Server Data Access}

A storage server operator (or a host compromise) can read fragment files at rest directly from
disk. Although one server typically holds only one fragment per block, repeated observation
across many blocks still discloses user data patterns and potentially sensitive content.

\subsubsection{I.3: Fingerprint Analysis}

An attacker who collects cross-checksum material (hashes, fingerprints, and metadata) can run
offline analysis to correlate uploads, infer duplicate content patterns, or support dictionary
style guessing against predictable objects. While fingerprints are not raw data, metadata leakage
still increases attacker knowledge about stored content.

\subsubsection{I.4: Error Message Information Leakage}

Verbose API errors and debug logs may expose internal identifiers, verification outcomes,
fragment indices, and topology details (server IDs and ports). This information can be used by
attackers to target specific nodes, optimize denial-of-service behavior, or refine tampering
attempts.

\subsection{Denial of Service}

% COMMENT: Threats where attackers disrupt legitimate system operation. Examples: exhausting
% storage capacity, flooding network with requests, corrupting fragments to make data unrecoverable,
% causing expensive cryptographic operations. Consider both resource exhaustion and functional
% denial of service.

\subsubsection{D.1: Storage Exhaustion}

% THREAT: An attacker stores excessive data, filling storage servers and preventing legitimate
% users from storing data.

\subsubsection{D.2: Fragment Unavailability}

% THREAT: An attacker compromises more than n-m servers or corrupts more than n-m fragments,
% making data unrecoverable.

\subsubsection{D.3: Verification Flooding}

% THREAT: An attacker triggers excessive verification operations, consuming CPU resources with
% expensive finite field arithmetic and cryptographic operations.

\subsubsection{D.4: Network Bandwidth Exhaustion}

% THREAT: An attacker floods the network with fragment requests or verification messages,
% preventing legitimate communication.

\subsection{Elevation of Privilege}

% COMMENT: Threats where attackers gain capabilities they are not authorized to have. Examples:
% exploiting buffer overflows to execute arbitrary code, bypassing access controls, escalating
% from client to server privileges. Consider both implementation vulnerabilities and design flaws.

\subsubsection{E.1: Code Injection via Malformed Input}

% THREAT: An attacker sends maliciously crafted fragments or fingerprints that exploit parsing
% vulnerabilities, leading to arbitrary code execution on servers.

\subsubsection{E.2: Verification Bypass}

% THREAT: An attacker discovers a way to bypass fingerprint verification, allowing corrupted
% fragments to be accepted as valid.

\subsubsection{E.3: Cryptographic Parameter Override}

% THREAT: An attacker manipulates system parameters (fingerprint size, field size, collision
% probability) to weaken security and enable attacks.

\subsubsection{E.4: Server-to-Server Privilege Escalation}

% THREAT: A compromised storage server exploits inter-server protocols to gain control over
% other servers or access fragments it should not have.

\newpage
\end{document}